% Corrected French for printed p.138.
% Corrections: D -> C (middle term); classe abéliennne -> catégorie
% abélienne; qussi-égal -> quasi-égal.
\SourcePage{138}

\LayerTarget{tohoku:def:subcategory:p2}
\noindent plus des conditions précédentes, les groupes
$\mathrm{Hom}_{\mathrm{C}'}(A,B)$ sont des sous-groupes de
$\mathrm{Hom}_{\mathrm{C}}(A,B)$.
\LayerTarget{tohoku:def:complete-subcategory}
Supposons que $\mathbf{C}$ soit une catégorie abélienne, $\mathbf{C}'$ est appelé une
sous-catégorie \textit{complète} si (i) \textit{pour $A,B\in \mathbf{C}'$, on a
$\mathrm{Hom}_{\mathrm{C}'}(A,B)=\mathrm{Hom}_{\mathrm{C}}(A,B)$.} (ii) \textit{Si dans une
suite exacte $A\to B\to C\to
D\to E$ les 4 termes extrêmes sont dans $\mathbf{C}'$, il en est de même
du terme médian $C$.} D'après (i), $\mathbf{C}'$ est complètement déterminé par la
classe de ses objets. (ii) équivaut à dire que pour tout morphisme
$P\to Q$ avec $P,Q\in \mathbf{C}'$ le noyau et le conoyau sont dans $\mathbf{C}'$,
et que pour toute suite exacte
$0\to R'\to R\to R''\to 0$
avec $R',R''$ dans $\mathbf{C}'$, $R$ est dans $\mathbf{C}'$. On vérifie aussitôt, qu'alors
$\mathbf{C}'$ est elle-même une classe abélienne, et que pour un morphisme
$u:A\to B$ dans $\mathbf{C}'$, noyau et conoyau (donc aussi image et
coimage) sont identiques au noyau et conoyau (resp. \ldots\,.) envisagé dans
$\mathbf{C}$.

\LayerTarget{tohoku:def:thick-subcategory}
La sous-catégorie $\mathbf{C}'$ de $\mathbf{C}$ est dite \textit{épaisse}, si elle satisfait à
la condition (i) ci-dessus, et au renforcement suivant de la condition (ii) :
(iii) \textit{Si dans une suite exacte
$A\to B\to C$ les termes extrêmes $A,C$ sont dans
$\mathbf{C}'$, il en est de même de $B$.} Si $\mathbf{C}$ est la catégorie abélienne des
groupes abéliens, on trouve la notion de ``classe de groupes abéliens'' de
[17]. On voit comme dans [17] que (iii) équivaut à la conjonction des trois
conditions suivantes : Tout objet nul est dans $\mathbf{C}'$, tout objet isomorphe à
un sous-truc ou un truc-quotient d'un objet de $\mathbf{C}'$ est dans $\mathbf{C}'$, toute
extension de deux objets de $\mathbf{C}'$ est dans $\mathbf{C}'$.

\LayerTarget{tohoku:def:quotient-category}
Soit $\mathbf{C}$ une catégorie abélienne, $\mathbf{C}'$ une sous-catégorie épaisse, nous allons
définir une nouvelle catégorie abélienne, notée $\mathbf{C}/\mathbf{C}'$ et appelée
\textit{catégorie quotient} de $\mathbf{C}$ par $\mathbf{C}'$. Les objets de $\mathbf{C}/\mathbf{C}'$ sont par
définition les objets de $\mathbf{C}$. Nous allons définir les morphismes dans $\mathbf{C}/\mathbf{C}'$
de $A$ dans $B$, appelés ``morphismes mod. $\mathbf{C}'$ de $A$ dans $B$''.
\LayerTarget{tohoku:def:quasi-equal}
Disons qu'un sous-truc $A'$ de $A$ est égal mod. $\mathbf{C}'$ à $A$, ou quasi-égal à
$A$, si $A/A'\in \mathbf{C}'$ ; alors tout sous-truc de $A$ contenant $A'$ est encore
quasi-égal à $A$, de plus le inf de deux sous-trucs de $A$ quasi-égaux à $A$
est encore quasi-égal à $A$. Dualement, on introduit la notion de quotient de
$B$ quasi-égal à $B$ : un tel quotient $B/N$ est quasi-égal à $B$ si
$N\in \mathbf{C}'$.
\LayerTarget{tohoku:def:quotient-morphism}
Un morphisme mod. $\mathbf{C}'$ de $A$ dans $B$ est alors défini par un morphisme $f'$
d'un sous-truc $A'$ de $A$ quasi-égal à $A$, dans un quotient $B'$ de $B$
quasi-égal à $B$, étant entendu qu'un morphisme
$f'':A''\to B''$ (satisfaisant aux mêmes conditions), définit le
même morphisme mod. $\mathbf{C}'$, si et seulement si on peut trouver
$A'''\leqq\mathrm{Inf}(A',A'')$, $B'''\leqq\mathrm{Inf}(B',B'')$, $A'''$ quasi
égal à $A$, $B'''$ quasi égal à $B$, tels que les morphismes
$A'''\to B'''$ définis par $f',f''$ soient identiques. Cette
dernière relation entre $f'$ et $f''$ est bien une relation d'équivalence, la
définition précédente des morphismes mod. $\mathbf{C}'$ est donc cohérente.
\LayerTarget{tohoku:def:quotient-hom}
Supposons que pour tout $A\in \mathbf{C}$, les sous-trucs de $A$ forment un ensemble
(ce qui est vrai pour toutes les catégories connues), alors on peut
considérer l'ensemble des morphismes mod. $\mathbf{C}'$ de $A$ dans $B$, noté
$\mathrm{Hom}_{\mathrm{C}/\mathrm{C}'}(A,B)$. $\mathrm{Hom}_{\mathrm{C}/\mathrm{C}'}(A,B)$ apparaît comme limite
inductive des groupes abéliens $\mathrm{Hom}_{\mathrm{C}}(A',B')$ où $A'$ ($B'$)
parcourt les sous-trucs de $A$ (quotients de $B$) quasi-égaux à $A$ ($B$), et
peut être considéré par suite comme un groupe abélien.
\LayerTarget{tohoku:def:quotient-composition:p1}
On définit de même un accouplement
$\mathrm{Hom}_{\mathrm{C}/\mathrm{C}'}(A,B)\times\mathrm{Hom}_{\mathrm{C}/\mathrm{C}'}(B,C)  \to\mathrm{Hom}_{\mathrm{C}/\mathrm{C}'}(A,C)$ de la façon suivante. Soit $u\in\mathrm{Hom}_{\mathrm{C}}(A',B')$ un
représentant d'un $u'\in\mathrm{Hom}_{\mathrm{C}/\mathrm{C}'}(A,B)$, et
$v\in\mathrm{Hom}_{\mathrm{C}}(B'',C')$ un représentant d'un
$v'\in\mathrm{Hom}_{\mathrm{C}/\mathrm{C}'}(B,C)$. Soit $Q$ l'image du morphisme canonique
$B''\to B'$ du

\TohokuContinuation{139}
